\begin{abstract}
Thi trực tuyến đặt ra nhu cầu kết hợp quản lý đề thi, chấm điểm và hỗ trợ giám sát trong cùng một quy trình. Đồ án EduShield AI xây dựng ứng dụng web cho sinh viên và giảng viên: nhập đề từ Excel hoặc Word, công bố đề, làm bài trắc nghiệm, chấm điểm tự động và tra cứu sự kiện kèm ảnh minh chứng.

Hệ thống sử dụng React ở giao diện, FastAPI ở máy chủ và SQLite thông qua SQLAlchemy để lưu dữ liệu. Xử lý hình ảnh được thực hiện trong trình duyệt bằng mô hình Face Landmarker và EfficientDet Lite0 đã được huấn luyện sẵn. Kết quả phát hiện khuôn mặt, điện thoại và các quy tắc hình học, độ sáng được kết hợp để sinh cảnh báo. Ảnh thu nhỏ chỉ được gửi khi phát sinh sự kiện nghi vấn, thay vì truyền liên tục toàn bộ video.

Báo cáo phân tích phiên bản mã nguồn \texttt{97a11fe} và thực hiện 17 ca kiểm thử chức năng trên dữ liệu tổng hợp; tất cả các ca trong bộ kiểm thử đều đạt. Bản cập nhật bổ sung 11 ca API về thời gian và sự kiện, trong đó hai ca tái hiện giới hạn chưa cưỡng chế hạn nộp. Các kết quả xác nhận một số hành vi của bộ nhập đề, API, cơ chế chấm điểm và kiểm soát truy cập. Độ chính xác nhận diện, độ trễ thực tế và khả năng chịu tải chưa được đo trong phạm vi đánh giá này. Các giới hạn đáng chú ý gồm định danh sinh viên tự khai, API nộp chưa cưỡng chế deadline, liên kết sự kiện với bài thi chưa trực tiếp và chưa có hàng đợi gửi lại bền vững. Hệ thống phù hợp làm nguyên mẫu hỗ trợ giám sát; quyết định xử lý vi phạm cần giảng viên xem xét ngữ cảnh và bằng chứng.

\noindent\textbf{Từ khóa:} thi trực tuyến; MediaPipe; giám sát hỗ trợ; FastAPI; React; nhận diện sự kiện.
\end{abstract}

Hệ thống mở rộng truy cập trên PC, Android và iPhone thông qua HTTPS do Cloudflare Tunnel cấp. FastAPI phục vụ cả giao diện build và API; thời điểm bắt đầu được lưu trong bảng \src{exam_starts}. Sự kiện chuyển tab, mất focus và thoát toàn màn hình được gửi về backend trong lúc làm bài. Android APK đã hoạt động theo xác nhận của tác giả; iOS sử dụng Safari/PWA, không được mô tả như một gói IPA native.
